\clearpage
\section{问题四：混合干扰源的联合搜索、定位与清除}

问题四在问题三的基础上加入发射方向未知的定向干扰源。设频道 $c$
的源位置为 $x_c$、有效接收半径为 $R_c\in[1000,1500]$；若该源为
定向源，其未知发射方向单位向量记为 $u_c$，则接收点 $p$ 能够收到
信号的条件为
\begin{equation}
  \|p-x_c\|_2\le R_c,\qquad
  (p-x_c)^{\mathsf T}u_c\ge0.
  \label{eq:p4-visibility}
\end{equation}
第二个不等式表示定向源只覆盖发射方向前方的闭半平面。因此，单点
无信号可能由频道空缺、距离超限或接收点位于辐射背面引起，不能沿用
问题三的单圆排除规则。本问仍对每个频道维护位置可行集 $\Omega_c$，
但将多方位搜索覆盖和多点无信号排除作为两个独立的几何约束。

\safeimage[0.99\textwidth]{figures/q4_method_overview.pdf}{问题四的混合源处理模型。（a）定向源只在有效接收半径内的前向半平面产生可接收信号，单点无信号不能直接排除源位置。（b）中心、内环和外环巡检点形成覆盖目标圆域的多方位接收布局。（c）多个无信号观测点共同确定保守排除区域，剩余位置可行集用于选择异方位复测点\label{fig:q4-framework}}

\subsection{混合源的集合成员观测模型}

记频道 $c$ 的有信号和无信号观测点集分别为
$\mathcal P_c^+$ 与 $\mathcal P_c^-$。获得示向度 $\alpha$ 时，沿用
问题三的有界误差示向扇形 $W(p,\alpha)$；返回近距离信号时，使用
$5\,\mathrm m$ 邻域约束。相应集合递推为
\begin{equation}
  \Omega_c\leftarrow
  \begin{cases}
    \Omega_c\cap W(p,\alpha)\cap B(p,1500),
      & \text{获得示向度},\\
    \Omega_c\cap B(p,5),
      & \text{返回近距离信号},\\
    \Omega_c\setminus B(p,20),
      & \text{在 $p$ 清除失败}.
  \end{cases}
  \label{eq:p4-positive-update}
\end{equation}
清除结果只取决于源与清除点的距离，故式~\eqref{eq:p4-positive-update}
中的光学排除对全向源和定向源均成立。对于单次无信号反馈，程序仅将
$p$ 加入 $\mathcal P_c^-$，不立即改变 $\Omega_c$。已获得有信号观测
后，全向源和定向源采用相同的最小二乘交会参考位置 $\hat x_c$ 与
最小覆盖圆半径 $r_c$，从而将方向不确定性限制在搜索和无信号解释
两个环节。

\subsection{多方位可探测覆盖条件}

设已经对全部未发现频道实施检测的巡检点集合为 $\mathcal S$。对
$\mathcal S$ 中任意三个不共线点 $p_i,p_j,p_k$，定义局部有效区域
\begin{equation}
  D_{ijk}=
  \operatorname{conv}\{p_i,p_j,p_k\}
  \cap B(p_i,1000)\cap B(p_j,1000)\cap B(p_k,1000),
  \label{eq:p4-local-region}
\end{equation}
并令
\begin{equation}
  \mathcal D_4(\mathcal S)
  =\bigcup_{\substack{i<j<k\\
  \max_{a,b\in\{i,j,k\}}\|p_a-p_b\|_2<2000}}
  D_{ijk}.
  \label{eq:p4-local-hull}
\end{equation}
对任意 $x\in D_{ijk}$，三个巡检点到 $x$ 的距离均不超过最小接收
半径 $1000\,\mathrm m$。同时，$x$ 位于三个点的凸包内；任意经过
$x$ 的闭半平面至少包含其中一个巡检点，否则三个点将全部位于另一
开半平面，与 $x$ 属于其凸包矛盾。由式~\eqref{eq:p4-visibility}，
无论源的发射方向如何，位于 $D_{ijk}$ 的定向源均能被至少一个巡检点
检测到。因此
\begin{equation}
  \mathcal T\subseteq\mathcal D_4(\mathcal S)
  \label{eq:p4-discovery-condition}
\end{equation}
是目标圆域内所有未清除源均已被搜索到的充分条件。当三角形最长边
不超过 $1000\,\mathrm m$ 时，式~\eqref{eq:p4-local-region} 中三个
圆盘约束自动成立；实现仍统一执行集合运算，以保持边界处理一致。

程序采用中心、内环和外环组成的固定巡检布局：中心点为原点，内环
含 $12$ 个半径 $925\,\mathrm m$ 的等角节点并相对 $x$ 轴偏转
$10\degree$，外环含 $18$ 个半径 $1850\,\mathrm m$ 的等角节点。
程序启动时直接以连续多边形运算验证式~\eqref{eq:p4-discovery-condition}；
若布局不能覆盖整个目标圆域则拒绝运行。覆盖状态只由已经实际完成
全频道检测的节点更新，规划中尚未访问的节点和局部定位点均不提前
计入 $\mathcal S$。

\subsection{多点无信号排除与异方位复测}

若已发现频道 $c$ 在至少三个彼此相距不小于 $1\,\mathrm m$ 的位置
得到无信号反馈，则由 $\mathcal P_c^-$ 按
式~\eqref{eq:p4-local-hull} 构造 $\mathcal D_4(\mathcal P_c^-)$。
若源位于该区域，则任意可能的发射半平面内都存在一个距离不超过
$1000\,\mathrm m$ 的无信号观测点，这与观测结果矛盾，故可作保守
更新
\begin{equation}
  \Omega_c\leftarrow
  \Omega_c\setminus
  \operatorname{int}_{\delta}
  \mathcal D_4(\mathcal P_c^-),
  \label{eq:p4-negative-update}
\end{equation}
其中 $\operatorname{int}_{\delta}$ 表示向内收缩一个数值容差后的
内部。该更新不会把单点背向失联误判为位置排除；若更新后
$\Omega_c$ 为空，则程序将其判定为观测矛盾并中止。

首次测向后，程序沿示向线在可行距离区间前移 $26.18\%$，并在法向
两侧设置 $82.76\,\mathrm m$ 的候选点；后续复测围绕 $\hat x_c$
构造法向候选点。当频道同时具有有信号和无信号历史时，将候选点关于
最近一次有信号观测的估计轴作镜像，选择到既有无信号点最远的一侧，
以增加再次进入辐射正面的机会。若当前位置对其他已发现频道满足
基线不小于 $158.29\,\mathrm m$、交会角正弦不小于 $0.15$ 且距离
小于 $1248.26\,\mathrm m$，则在原地执行跨频道补充测向。

\subsection{滚动调度、清除与终止准则}

每个决策时刻将未访问的计划巡检点与已发现频道的局部动作合并为
候选任务集。程序以最近邻开放路径为初解，执行 $2000$ 次二边交换
模拟退火并保留最短路径；与问题三相同，每轮只执行首个任务，获得
新观测后更新 $\Omega_c$、$\mathcal S$ 和候选任务集，再重新规划。
局部定位阶段不实施额外的全频道巡检，避免把未按全频道检测的局部
停驻点错误计入式~\eqref{eq:p4-discovery-condition}。

当 $r_c\le19.95\,\mathrm m$ 时，程序选择能够使
$\Omega_c\subseteq B(q,20)$ 且距当前位置最近的清除点 $q$。
当 $19.95<r_c\le65.71\,\mathrm m$ 时，允许在 $\hat x_c$ 处进行
至多两次试探性清除，失败后按式~\eqref{eq:p4-positive-update}
继续收缩可行集。每个频道的常规局部动作上限为 $10$ 次；达到上限
后，以边长 $25\,\mathrm m$ 的方格中心覆盖剩余可行集。方格半对角线
$25\sqrt2/2=17.68\,\mathrm m<20\,\mathrm m$，故该有限网格能够提供
确定性的光学清除兜底。

记 $\mathcal K_k$ 为时刻 $k$ 已发现但尚未清除的频道集合，
$N_k^{\mathrm{clr}}$ 为累计清除数。控制器的终止准则为
\begin{equation}
  N_k^{\mathrm{clr}}=16
  \quad\lor\quad
  \left(\mathcal K_k=\varnothing
  \ \land\
  \mathcal T\subseteq\mathcal D_4(\mathcal S_k)\right).
  \label{eq:p4-stop}
\end{equation}
第一项利用题设给出的源数量上界；源数量不足 $16$ 时，第二项保证
所有可能位置和发射方向均已完成搜索，且已发现源均已清除。全局动作
上限为 $5500$，同时按 \texttt{/enter} 返回的现实剩余时间设置截止
约束。

\begin{modelalgorithm}{混合干扰源的多方位搜索与保守清除}
  \label{alg:p4}
  \AlgoIn{目标圆域、频道集合、接收距离范围和多方位巡检参数}
  \AlgoOut{清除记录、逐频道位置可行集和行为日志}
  \begin{enumerate}[leftmargin=2.5em]
    \AlgoStep{生成中心、内环和外环巡检节点，并验证式~\eqref{eq:p4-discovery-condition}。}
    \AlgoStep{调用 \texttt{/enter}，在原点检测全部频道，初始化 $\mathcal P_c^+$、$\mathcal P_c^-$ 与 $\Omega_c$。}
    \AlgoStep{合并未执行巡检节点和已发现频道的局部任务，优化开放路径并执行首个动作。}
    \AlgoStep{按式~\eqref{eq:p4-positive-update}更新有信号观测；满足多点条件时按式~\eqref{eq:p4-negative-update}排除区域。}
    \AlgoStep{依据 $r_c$ 选择异方位复测、集合包含清除或有限网格光学覆盖。}
    \AlgoStep{满足式~\eqref{eq:p4-stop}后调用 \texttt{/exit}，否则返回步骤 3。}
  \end{enumerate}
\end{modelalgorithm}

\subsection{分层随机实验与极端构型检验}

为检验方向不确定性增加后的稳定性，本文设置三类逐级加严的固定种子
随机实验，每类运行 $500$ 次且均含 $16$ 个源。混合源场景包含
$8$ 个全向源和 $8$ 个随机朝向定向源；全定向随机场景的 $16$ 个源
均为随机位置、随机朝向定向源；径向外射场景进一步将源限制在半径
$1150$--$1800\,\mathrm m$ 的环带，并令发射方向沿位置矢量向外。
第三类场景使圆心和多数内环点位于信号背面，用于集中检验外环节点的
多方位接收能力。各场景的频道、接收半径和示向误差仍按题设范围独立
生成，控制器参数在全部运行中保持不变。

评价指标沿用问题三：清除比例为 $\rho=N_c/N$，完整清除样例数为
$N_c=N$ 的运行次数；总虚拟时间的均值、标准差和 $95\%$ 分位数仅对
完整运行计算，避免中途异常产生的部分时间压低统计结果。

\begin{table}[H]
  \centering
  \caption{问题四三类分层随机实验结果}
  \label{tab:p4-monte-carlo}
  \small
  \begin{tabularx}{\textwidth}{
    >{\centering\arraybackslash}p{0.22\textwidth}
    >{\centering\arraybackslash}p{0.08\textwidth}
    >{\centering\arraybackslash}p{0.14\textwidth}
    *{2}{>{\centering\arraybackslash}X}}
    \toprule
    场景 & 样例数 & 完整清除样例数
    & 总虚拟时间均值与标准差/s & 总虚拟时间95\%分位数/s \\
    \midrule
    混合源随机朝向 & 500 & 499
    & $5577.2\pm777.8$ & 6636.1 \\
    全定向源随机朝向 & 500 & 500
    & $6443.1\pm687.0$ & 7493.6 \\
    全定向源径向外射 & 500 & 500
    & $7857.0\pm396.8$ & 8542.9 \\
    \bottomrule
  \end{tabularx}
\end{table}

\safeimage[0.99\textwidth]{figures/q4_validation_overview.pdf}{问题四的极端构型与时间分布。（a）一个全定向源径向外射样例，短箭头表示发射方向，灰线连接干扰源与能够接收其信号的最近计划巡检点；该样例的 $16$ 个源均至少对应一个有效巡检点。（b）三类场景各 $500$ 次运行的总虚拟时间分布，白色箱体表示四分位区间和中位数\label{fig:q4-validation}}

全部 $1500$ 个样例中有 $1499$ 个完成 $16/16$ 清除，样本内完整
清除率为 $99.93\%$，其 Wilson $95\%$ 置信下界为 $99.62\%$。
唯一未完成的混合源样例在清除 $11$ 个源后因多边形布尔运算触发
拓扑异常而中止，并非达到动作数或虚拟时间上限。表中时间统计仅使用
完整运行：总体均值为 $6626.4\,\mathrm s$，$95\%$ 分位数为
$8196.8\,\mathrm s$，最大值为 $9533.9\,\mathrm s$。

图~\ref{fig:q4-validation}(a) 中，尽管所有定向源均背离圆心发射，
每个源仍至少对应一个位于其前向半平面且距离不超过有效接收半径的
外环巡检点。图~\ref{fig:q4-validation}(b) 表明，定向源比例和朝向
限制增强时，总时间分布整体右移；混合源、全定向随机和径向外射场景
的平均时间依次增加。该结果说明主要新增代价来自背向失联后的异方位
复测，而中心--双环布局在径向外射压力场景中仍保持了完整搜索能力。

\subsection{正式测试结果}

三次正式测试用于记录控制器在竞赛模拟器中的实际执行结果和接口运行
情况，不用于替代前述 $1500$ 次随机实验对模型总体表现的评价。平均
定位清除时间按模拟器返回的总虚拟时间除以成功清除数量计算，程序运行
时间采用实际墙钟时间。

\begin{table}[H]
  \centering
  \caption{问题四三次正式测试结果}
  \label{tab:p4-official}
  \begin{tabular}{ccccc}
    \toprule
    测试 & 测试案例编码 & 清除干扰源个数
    & 平均定位清除时间/s & 程序运行时间/s \\
    \midrule
    测试 1 & \tableplaceholder & \tableplaceholder & \tableplaceholder & \tableplaceholder \\
    测试 2 & \tableplaceholder & \tableplaceholder & \tableplaceholder & \tableplaceholder \\
    测试 3 & \tableplaceholder & \tableplaceholder & \tableplaceholder & \tableplaceholder \\
    \bottomrule
  \end{tabular}
\end{table}
